% Part: incompleteness
% Chapter: incompleteness-provability
% Section: godels-paper

\documentclass[../../../include/open-logic-section]{subfiles}

\begin{document}

\olfileid[id]{inc}{inp}{gop}

\olsection{Perbandingan dengan Makalah Asli G\"odel}

Ada baiknya meluangkan waktu untuk mempelajari makalah G\"odel tahun 1931.
Pengantarnya menguraikan secara garis besar gagasan-gagasan yang baru saja
kita bahas. Bahkan jika Anda hanya membaca sekilas makalah itu, mudah untuk
melihat apa yang terjadi pada setiap tahap: mula-mula G\"odel menguraikan
sistem formal $P$ (sintaksis, aksioma, aturan pembuktian); kemudian ia
mendefinisikan fungsi dan relasi rekursif primitif; lalu ia menunjukkan bahwa
$x B y$ rekursif primitif, dan berargumen bahwa fungsi dan relasi rekursif
primitif direpresentasikan dalam~$\Th{P}$. Setelah itu ia membuktikan teorema
ketaklengkapan seperti di atas. Dalam Bagian~3, ia menunjukkan bahwa kita
dapat memilih pernyataan yang tidak dapat dibuktikan sebagai suatu kalimat
dalam bahasa aritmetika. Inilah asal-usul lema~$\beta$, yang juga kita
gunakan untuk menangani barisan ketika menunjukkan bahwa fungsi rekursif
dapat direpresentasikan dalam~$\Th{Q}$. G\"odel tidak sampai mengisolasi
suatu himpunan aksioma minimal yang memadai, tetapi sekarang kita mengetahui
bahwa $\Th{Q}$ sudah mencukupi. Akhirnya, dalam Bagian~4, ia membuat sketsa
pembuktian teorema ketaklengkapan kedua.

\end{document}
